\chapter{Ricarica intelligente di veicoli elettrici}
\label{cap:ricaricaev}

\noindent\textbf{Classe:} LP / QP convesso \hfill
\textbf{Script:} \texttt{python/lab10\_ricarica\_ev.py}

\section{Motivazione gestionale}

Una flotta aziendale deve arrivare carica al mattino sfruttando le ore notturne a
prezzo basso --- ma se tutti i veicoli caricano insieme nelle ore economiche, il
prelievo di punta esplode e con esso i costi di potenza e il rischio di superare la
capacità del contatore. Il modello decide \emph{quanta potenza} dare a ciascun veicolo
in ciascuna ora: una decisione naturalmente continua, senza alcun bisogno di variabili
di accensione/spegnimento.



La ricarica intelligente (\emph{smart charging}) è uno dei problemi operativi nuovi
creati dalla transizione energetica: flotte elettriche, prezzi orari dell'energia,
limiti di potenza dei contatori. È anche un esempio perfetto di problema che
\emph{sembra} richiedere variabili di accensione/spegnimento e invece è un LP puro,
perché la decisione vera --- quanta potenza --- è continua.

\textbf{In questo capitolo impareremo a:} (1) modellare decisioni su un orizzonte
orario con finestre di disponibilità; (2) leggere i duali come ``prezzi delle ore
marginali''; (3) gestire il trade-off costo-picco con un approccio multiobiettivo;
(4) diagnosticare un'insidia reale dell'analisi di sensitività (il RHS con
costanti).

\section{Costruzione guidata del modello}

\paragraph{Il problema a parole.} \emph{Decidiamo} quanta potenza erogare a ogni
veicolo in ogni ora. \emph{L'obiettivo} è minimizzare la spesa energetica della
notte. \emph{I vincoli}: ogni veicolo deve ricevere tutta l'energia che gli serve
entro la partenza, può caricare solo quando è collegato e mai oltre la potenza del
suo caricatore, e il prelievo complessivo (ricarica più edificio) non può mai
superare la potenza del contatore.

\subsection*{Insiemi e dati (input del modello)}
\begin{center}\small
\begin{tabular}{llp{7.2cm}}
\toprule
Simbolo & Tipo & Significato \\
\midrule
$V$, $T$ & insiemi finiti & veicoli e ore (qui $\Delta t = 1$ h) \\
$\pi_t$ & $\in \R_{\ge 0}$ & prezzo dell'energia nell'ora $t \in T$ (\euro/kWh) \\
$a_{vt}$ & $\in \{0, 1\}$ & \emph{dato} osservato: 1 se il veicolo $v$ è collegato nell'ora $t$ \\
$e_v$ & $\in \R_{> 0}$ & energia richiesta dal veicolo $v$ (kWh) \\
$\bar p_v$ & $\in \R_{> 0}$ & potenza massima di ricarica del veicolo $v$ (kW) \\
$b_t$ & $\in \R_{\ge 0}$ & carico di base dell'edificio nell'ora $t$ (kW) \\
$k$ & $\in \R_{> 0}$ & potenza massima prelevabile dal contatore (kW) \\
$\eta$ & $\in (0, 1]$ & rendimento di carica \\
\bottomrule
\end{tabular}
\end{center}

\subsection*{Variabili decisionali}
\begin{center}\small
\begin{tabular}{llp{7.2cm}}
\toprule
$x_{vt}$ & $\ge 0$, continua & potenza assegnata al veicolo $v \in V$ nell'ora $t \in T$ (kW) \\
\bottomrule
\end{tabular}
\end{center}

\begin{modello}[Ricarica a costo minimo]
\begin{align}
\min\;& \sum_{v \in V}\sum_{t \in T} \pi_t\, \Delta t\, x_{vt} \notag\\
\text{s.t.}\;\;
& \eta \sum_{t \in T} \Delta t\, x_{vt} \ge e_v && \forall v \in V \label{eq:energia}\\
& 0 \le x_{vt} \le a_{vt}\, \bar p_v && \forall v \in V,\ t \in T \label{eq:potenza}\\
& \sum_{v \in V} x_{vt} + b_t \le k && \forall t \in T \label{eq:rete}
\end{align}
Varianti: peak shaving ($\min z$ con $\sum_{v \in V} x_{vt} + b_t \le z$ per ogni
$t \in T$); profilo regolare (QP con $\gamma \sum_{t \in T,\, t > 1} (z_t - z_{t-1})^2$);
multiobiettivo costo $+\,\rho\,z$.
\end{modello}

\begin{spiegazione}
\textbf{Funzione obiettivo.} La spesa energetica totale: la potenza $x_{vt}$ erogata
per un'ora ($\Delta t$) costa $\pi_t\, \Delta t\, x_{vt}$ euro; sommiamo su veicoli
e ore.

\textbf{Perché niente variabili binarie.} La disponibilità $a_{vt}$ è un \emph{dato}
(il veicolo è in deposito oppure no), non una decisione: entra nel bound
\eqref{eq:potenza} spegnendo le ore in cui il veicolo è assente. La decisione ---
quanta potenza erogare --- è naturalmente continua.

\textbf{Vincolo \eqref{eq:energia}.} Scritto come $\ge$: caricare più del necessario è
ammesso ma mai conveniente (il prezzo è positivo), quindi all'ottimo vale con
uguaglianza. Il rendimento $\eta$ sta a sinistra: dell'energia prelevata solo il 95\%
finisce in batteria.

\textbf{Vincolo \eqref{eq:rete}.} Accoppia i veicoli tra loro: senza di esso il
problema si separerebbe per veicolo (ognuno caricherebbe nelle proprie ore più
economiche).
\end{spiegazione}

\section{Esempio svolto a mano}

\begin{esempio}[Un veicolo, due ore]
Servono $e = 10$ kWh ($\eta = 1$); prezzi $\pi_1 = 0{,}10$ e $\pi_2 = 0{,}20$
\euro/kWh; potenza massima $\bar p = 8$ kW.

\textbf{Caso $\bar p = 8$.} Nell'ora economica si carica al massimo ($8$ kWh); i restanti
$2$ kWh vanno nell'ora cara. Costo: $8 \cdot 0{,}10 + 2 \cdot 0{,}20 = 1{,}20$~\euro.
\emph{Prezzo ombra del fabbisogno}: l'undicesimo kWh finirebbe nell'ora cara
$\Rightarrow$ $0{,}20$ \euro/kWh. \emph{Prezzo ombra del limite di potenza} nell'ora
1: un kW in più permetterebbe di spostare un kWh dall'ora cara a quella economica
$\Rightarrow$ $0{,}10$ \euro/kW.

\textbf{Caso $\bar p = 12$.} Tutto nell'ora economica: costo $1{,}00$~\euro, duale del
fabbisogno $0{,}10$, duale della potenza $0$ (vincolo non attivo). Il pattern
``il duale è il prezzo dell'ora marginale'' è la chiave di lettura di tutto il
capitolo.
\end{esempio}

\section{Caso di studio}

Sei furgoni con finestre di ricarica notturne diverse (dati in
\texttt{dati/ev\_flotta.csv}), prezzi orari stilizzati con picco serale, carico di
base dell'edificio, contatore da $k = 120$ kW, $\eta = 0{,}95$.

\begin{lstlisting}[style=python]
x = m.addVars(veicoli, ore, name="x")
for v in veicoli:
    for t in ore:
        x[v, t].UB = P[v] * disp[v, t]    # 0 se non collegato
v_ene  = m.addConstrs((eta * gp.quicksum(x[v, t] for t in ore)
                       >= E[v] for v in veicoli), name="energia")
v_rete = m.addConstrs((gp.quicksum(x[v, t] for v in veicoli)
                       + base[t] <= C for t in ore), name="rete")
m.setObjective(gp.quicksum(prezzo[t] * x[v, t]
                           for v in veicoli for t in ore), GRB.MINIMIZE)
\end{lstlisting}

\section{Risultati}

\begin{lstlisting}[style=output]
Costo minimo   : 20,36 EUR/notte   picco di prelievo 103,4 kW (limite 120)
Peak shaving   : picco minimo possibile 68,0 kW
Prezzi ombra del fabbisogno (EUR/kWh):
  V1 0,0947   V2 0,0842   V3 0,0842   V4 0,0947   V5 0,0947   V6 0,0842
\end{lstlisting}

\begin{center}
\begin{minipage}{\textwidth}
\centering
\begin{tikzpicture}
\begin{axis}[lab, width=0.94\textwidth, height=4.0cm, ybar, bar width=3.2pt,
  title={Prezzo orario dell'energia}, xlabel={}, ylabel={c\euro/kWh},
  xmin=-0.7, xmax=23.7, ymin=0]
\addplot [fill=grigiomedio!60, draw=none] table [col sep=comma, x=ora, y=prezzo_cent]
  {\figdat{cap10_profili}};
\end{axis}
\end{tikzpicture}

\begin{tikzpicture}
\begin{axis}[lab, width=0.94\textwidth, height=6.0cm, const plot,
  title={Prelievo totale: inseguire i prezzi crea un picco notturno},
  xlabel={ora del giorno}, ylabel={prelievo (kW)}, xmin=-0.7, xmax=23.7, ymin=0, ymax=130]
\addplot [grigiomedio, dashed, thick] table [col sep=comma, x=ora, y=base]
  {\figdat{cap10_profili}};
\addplot [teal, very thick] table [col sep=comma, x=ora, y=totcosto]
  {\figdat{cap10_profili}};
\addplot [arancio, very thick] table [col sep=comma, x=ora, y=totpicco]
  {\figdat{cap10_profili}};
\addplot [rossomattone, densely dotted, thick] coordinates {(-0.7,120) (23.7,120)};
\legend{carico di base, costo minimo, peak shaving, limite rete}
\end{axis}
\end{tikzpicture}
\captionof{figure}{Sopra: prezzo orario dell'energia (barre grigie), basso di notte
e con picco serale. Sotto: prelievo totale (base $+$ ricarica) ora per ora nelle due
strategie: il costo minimo (teal) concentra la ricarica nelle ore economiche e crea
un picco notturno di 103 kW; il peak shaving (arancione) distribuisce la ricarica e
tiene il prelievo a 68 kW, sotto il limite di rete (linea rossa).}
\end{minipage}
\end{center}

\begin{spiegazione}[I duali raccontano le ore marginali]
I prezzi ombra del fabbisogno valgono $0{,}0842 = 0{,}08/\eta$ oppure
$0{,}0947 = 0{,}09/\eta$: sono i prezzi delle \emph{ore marginali} di ciascun veicolo
(l'ora più economica ancora libera nella sua finestra), corretti per il rendimento.
V2, V3 e V6 hanno finestre che coprono le ore a $0{,}08$; V1, V4 e V5 arrivano tardi
o partono presto e la loro ora marginale costa $0{,}09$. Un kWh di fabbisogno in più
costa quindi $8{,}4$--$9{,}5$ centesimi a seconda del veicolo.
\end{spiegazione}

\section{Analisi di sensitività}

\begin{lstlisting}[style=output]
Compromesso costo + rho*picco:
  rho = 0,00: costo 20,36  picco 103,4     rho = 0,20: costo 22,15  picco 68,0
  rho = 0,05: costo 20,69  picco  86,5     rho = 0,50: costo 22,15  picco 68,0
  rho = 0,10: costo 21,39  picco  74,9
Capacita' del contatore:
  k = 60, 65 kW: INAMMISSIBILE     k = 80: 21,09     k = 120: 20,36
  k = 70 kW    : 21,93             k = 90: 20,63
\end{lstlisting}

\begin{center}
\begin{tikzpicture}
\begin{axis}[lab, title={Frontiera costo-picco},
  xlabel={picco di prelievo (kW)}, ylabel={costo di ricarica (\euro)}]
\addplot+ [mark=*, mark size=1.8pt] table [col sep=comma, x=picco, y=costo]
  {\figdat{cap10_frontiera}};
\node [font=\scriptsize, anchor=west] at (axis cs:104,20.55) {$\rho=0$};
\node [font=\scriptsize, anchor=south] at (axis cs:68,22.3) {$\rho\ge 0{,}2$};
\end{axis}
\end{tikzpicture}
\captionof{figure}{Frontiera costo-picco ottenuta variando il peso $\rho$ del picco
nell'obiettivo: ogni punto è una strategia di ricarica ottima. Ridurre il picco da
103 a 75 kW costa circa un euro a notte; il punto $\rho \ge 0{,}2$ raggiunge il
picco minimo possibile (68 kW) al costo più basso che lo consente.}
\end{center}

\begin{insight}
Tagliare il picco da 103 a 75 kW ($-28$\%) costa solo un euro a notte; arrivare a
68 kW ne costa meno di due. Se la tariffa di potenza o l'upgrade del contatore
costano più di così, il peak shaving si ripaga da solo. Nota: il minimax puro
(picco 68) senza termine di costo spenderebbe 27,79~\euro{} --- il compromesso con
$\rho = 0{,}2$ ottiene \emph{lo stesso picco} a 22,15~\euro: mai ottimizzare un solo
obiettivo quando ce ne sono due.
\end{insight}

\begin{attenzione}[Un bug istruttivo, capitato davvero]
Nel vincolo \texttt{quicksum(x) + base[t] <= C}, Gurobi sposta la costante
\texttt{base[t]} nel termine noto: il RHS memorizzato è \texttt{C - base[t]}.
La prima versione della nostra analisi di sensitività scriveva
\texttt{v.RHS = nuovaC} e otteneva costi identici per ogni capacità --- stava in
realtà allentando il vincolo di \texttt{base[t]}. La forma corretta è
\texttt{v.RHS = nuovaC - base[t]}. Quando una sensitività non cambia nulla,
sospettare prima del proprio codice che del modello.
\end{attenzione}

\section{Esercizi}

\begin{esercizio}[verifica]
Verificare per perturbazione il prezzo ombra del fabbisogno di V1 (portare $e_1$ da
46 a 47 kWh) e spiegare perché vale $0{,}09/0{,}95$.
\end{esercizio}

\begin{esercizio}[arrivo in ritardo]
V3 arriva alle 23 invece che alle 20: come cambiano costo e picco? Quale prezzo
ombra aumenta?
\end{esercizio}

\begin{esercizio}[profilo regolare]
Implementare la variante QP con $\gamma \sum_{t \in T,\, t>1} (z_t - z_{t-1})^2$ dove
$z_t = b_t + \sum_{v \in V} x_{vt}$, per $\gamma \in \{0{,}01;\, 0{,}1;\, 1\}$: confrontare
costo, picco e ``nervosismo'' del profilo.
\end{esercizio}

\begin{esercizio}[emissioni]
Con intensità carbonica oraria $g_t$ (alta la sera, bassa di notte), tracciare la
frontiera costo-emissioni. Le ore economiche sono anche le più pulite in questo caso?
\end{esercizio}

\begin{esercizio}[capacità minima]
Trovare per bisezione la capacità minima $k$ che mantiene il problema ammissibile,
e spiegare quale combinazione di finestre e $\bar p_v$ la determina.
\end{esercizio}
